\scanpage{64}
seule fois, mais qu'il faille encore une construction transfinie.

Il y a donc des énoncés possibles pour caractériser les $R\in\operatorname{Ob}\Delta$ tels que $T_R$ soit une théorie \uline{algébrique}

1) $\exists$ \uline{petite} catégorie $J$, et ensembles $\Sigma$, \deleted{$\Sigma'$} de flèches de $\widetilde J$, \inserted{(\inserted{$+$ ens} $\Sigma'$ de \inserted{couples de} flèches de $\widetilde J$)} tels que $R$ se déduise de $\widetilde J$ en rendant inv. les flèches de $\Sigma$ (et égalisant les \inserted{couples de} flèches de $\Sigma'$)

2) \deleted{$\exists$ petite catégorie $J$, et ens} $\exists$ \deleted{famille} \inserted{pseudo-syst. inductif} transfini $(R_\alpha)$ dans $\Delta$, tel que
\[
 \begin{aligned}
 \text{(i)}\quad &R_0=\widetilde J,\quad J\text{ petite cat.}\\
 \text{(ii)}\quad &R_{\alpha+1}\text{ se déduit \unsure{en|de} $R_\alpha$ en rendant inversibles un}\\
 &\text{paquet de flèches, \uline{et} égalisant un autre paquet}\\
 &\text{de couples de flèches}\\
 \text{(iii)}\quad &R_\alpha=\varinjlim_{\beta<\alpha}R_\beta\quad\text{si $\alpha$ ordinal limite}\\
 \text{(iv)}\quad &R_\alpha=R\quad\text{pour $\alpha$ assez grand, $\alpha\in\mathcal U$ (l'Univers).}
 \end{aligned}
\]

\textcircled{2} \uline{Autre Pb :} caractériser les théories algébriques par des propriétés intrinsèques (plutôt que par ``construction'' via une suite transfinie de constructions, ou la représentabilité par un $R$ tel que\ldots).
La seule condition (probablement pas suffisante) que je voie est : \uline{$T$ commute aux 2-limites projectives de $\Delta$}, i.e. : $T$ transforme \deleted{\unsure{foncteurs}} \inserted{\unsure{ens|sans}} pl. fid (prod.) en \unsure{itou|isos|idem}, $\underline{\operatorname{Hom}}_{\mathrm{Cat}}(J,C)$ en $\operatorname{Hom}_{\mathrm{Cat}}(J,T(C))$\ldots

\textcircled{3} \uline{Conditions \deleted{\unsure{de $f$}} d'exactitudes généralisées} (comme suggérées par Schanuel) : on donne \deleted{\unsure{$R\to$}} $\rho\to\sigma$ dans $\mathcal R$ et on demande sur $C$ que $\underline{\operatorname{Hom}}_\Delta(\sigma,C)\to\underline{\operatorname{Hom}}_\Delta(\rho,C)$ soit une équivalence (resp. \inserted{pl.} fidèle, resp fidèle ?\ldots). Cela inclut \inserted{p. ex. rendre des flèches égales (pas seulement iso).} Les remarques sur l'existence de $\varprojlim$ et $\varinjlim$ dans $\Delta$ aussi restent mutat\unsure{.} mutand\unsure{.} applicables.
% archive sheet number: 63
